\subsection{两个复形的张量积}

设 (C, d) 为右 A-模复形，(C', d') 为左 A-模复形。

赋予k-模 \(C \otimes_{A} C'\) 以如下分次，使得

\[
(\mathrm{C} \otimes_ {\mathbf {A}} \mathrm{C} ^ {\prime}) _ {n} = \sum_ {p + q = n} (\mathrm{C} _ {p} \otimes \mathrm{C} _ {q} ^ {\prime})
\]

并令D表示唯一的次数为 \((-1)\) 的k-线性自同态，它作用于 C \(\otimes_{A} C'\)，使得

\[
\mathbf {D} (x \otimes x ^ {\prime}) = d x \otimes x ^ {\prime} + (- 1) ^ {p} x \otimes d ^ {\prime} x ^ {\prime}, \quad x \in \mathrm{C} _ {p}, y \in \mathrm{C} _ {q} ^ {\prime}, p, q \in \mathbf {Z}. \tag {1}
\]

我们有 \(D \circ D = 0\)，因为依（1）中的记号，

\[
\mathrm{D} ^ {2} (x \otimes x ^ {\prime}) = d d x \otimes x ^ {\prime} + (- 1) ^ {p - 1} d x \otimes d ^ {\prime} x ^ {\prime} + (- 1) ^ {p} d x \otimes d ^ {\prime} x ^ {\prime} - x \otimes d ^ {\prime} d ^ {\prime} x ^ {\prime}.
\]

\(k\) -模复形 \((\mathbf{C} \otimes_{\mathbf{A}} \mathbf{C}', \mathbf{D})\) 称为复形 \((\mathbf{C}, d)\) 与 \((\mathbf{C}', d')\) 的张量积复形。

注记. --- 1) 当 \(\mathbf{C}'\) 简化为 \(\mathbf{C}_0' = \mathbf{M}\) 时，则 \((\mathbf{C} \otimes_{\mathbf{A}} \mathbf{C}')_n = \mathbf{C}_n \otimes_{\mathbf{A}} \mathbf{M}\) 且 \(D = d \otimes 1_M\)；例如 \(\mathbf{C} \otimes_{\mathbf{A}} \mathbf{A}_s\) 自然等同于 \(C\) 。类似地，当 \(C\) 简化为 \(\mathbf{C}_0 = P\) 时，则 \((\mathbf{C} \otimes_{\mathbf{A}} \mathbf{C}')_n = P \otimes_{\mathbf{A}} \mathbf{C}_n'\) 且 \(D = 1_P \otimes d\)。

\begin{enumerate}
\def\labelenumi{\arabic{enumi})}
\setcounter{enumi}{1}
\tightlist
\item
  对每个整数 \(r\)，我们有 \((C \otimes_{A} C')(r) = C(r) \otimes_{A} C'\) ，但 \((C \otimes_{A} C')(r)\) 与 \(C \otimes_{A} C'(r)\) 一般不具有相同的微分。
\end{enumerate}

设 \(p, q\) 是两个整数， \(x \in Z_p(C)\) , \(x' \in Z_q(C')\) ；则元素 \(x \otimes x'\) 作为 \(C_p \otimes C_q'\) 的元素依(1)属于 \(Z_{p+q}(C \otimes C')\)；此外，若 \(y \in C_{p+1}\) , \(y' \in C_{q+1}'\) ，我们有

\[
(x + d y) \otimes (x ^ {\prime} + d ^ {\prime} y ^ {\prime}) = x \otimes x ^ {\prime} + \mathrm{D} (y \otimes x ^ {\prime} + (- 1) ^ {p} (x + d y) \otimes y ^ {\prime})
\]

通过取商，我们导出一个k-线性映射，称为典范映射

\[
\gamma_ {p, q} (\mathrm{C}, \mathrm{C} ^ {\prime}): \mathrm{H} _ {p} (\mathrm{C}) \otimes_ {\mathrm{A}} \mathrm{H} _ {q} (\mathrm{C} ^ {\prime}) \rightarrow \mathrm{H} _ {p + q} (\mathrm{C} \otimes_ {\mathrm{A}} \mathrm{C} ^ {\prime});
\]

如果我们装备 \(\mathbf{H}(\mathbf{C})\otimes_{\mathbf{A}}\mathbf{H}(\mathbf{C}^{\prime})\) 以如下分次，使得

\[
\left(\mathrm{H} (\mathrm{C}) \otimes \mathrm{H} (\mathrm{C} ^ {\prime})\right) _ {n} = \sum_ {p + q = n} \mathrm{H} _ {p} (\mathrm{C}) \otimes_ {\mathrm{A}} \mathrm{H} _ {p} (\mathrm{C} ^ {\prime}),
\]

\(\gamma_{p,q}\) 定义一个次数为0的分次 \(k\) -线性映射

\[
\gamma (\mathbf {C}, \mathbf {C} ^ {\prime}): \mathrm{H} (\mathbf {C}) \otimes_ {\mathbf {A}} \mathrm{H} (\mathbf {C} ^ {\prime}) \rightarrow \mathrm{H} (\mathbf {C} \otimes_ {\mathbf {A}} \mathbf {C} ^ {\prime}).
\]

II，第59页的命题6改述如下：

如果复形C和C'在右侧为零，则C ⊗\_A C' 在右侧为零，且典范的 k-线性映射

\[
\gamma_ {0, 0} (\mathbf {C}, \mathbf {C} ^ {\prime}): \mathrm{H} _ {0} (\mathbf {C}) \otimes_ {\mathbf {A}} \mathrm{H} _ {0} (\mathbf {C} ^ {\prime}) \rightarrow \mathrm{H} _ {0} (\mathbf {C} \otimes_ {\mathbf {A}} \mathbf {C} ^ {\prime})
\]

是双射。

设 \(u: (\mathbf{C}, d) \to (\mathbf{C}_1, d_1)\) 是右A-模复形的一个态射，且 \(u': (\mathbf{C}', d') \to (\mathbf{C}_1', d_1')\) 是左A-模复形的一个态射；则 \(u \otimes u': \mathbf{C} \otimes_{\mathbf{A}} \mathbf{C}' \to \mathbf{C}_1 \otimes_{\mathbf{A}} \mathbf{C}_1'\) 是 \(k\) -模复形的态射；事实上，它是次数为0的分次态射，并且若用D和 \(\mathbf{D}_1\) 表示 \(\mathbf{C} \otimes \mathbf{C}'\) 与 \(\mathbf{C}_1 \otimes \mathbf{C}_1'\) 的微分，则对于 \(p, q \in \mathbf{Z}, x \in \mathbf{C}_p, x' \in \mathbf{C}_q\) 我们有，

\[
\begin{array}{c} (u \otimes u ^ {\prime}) (\mathrm{D} (x \otimes x ^ {\prime})) = u (d x) \otimes u ^ {\prime} (x ^ {\prime}) + (- 1) ^ {p} u (x) \otimes u ^ {\prime} (d ^ {\prime} x ^ {\prime}) = \\ = d _ {1} u (x) \otimes u ^ {\prime} (x ^ {\prime}) + (- 1) ^ {p} u (x) \otimes d _ {1} ^ {\prime} u ^ {\prime} (x ^ {\prime}) = \mathrm{D} _ {1} (u (x) \otimes u ^ {\prime} (x ^ {\prime}))  . \end{array}
\]

此外，以下图表是可交换的：

\[
\begin{array}{c} \mathrm{H(C)} \otimes_ {\mathrm{A}} \mathrm {H(C^ {\prime})} \xrightarrow {\gamma (\mathrm{C,C} ^ {\prime})} \mathrm{H(C} \otimes_ {\mathrm{A}} \mathrm {C^ {\prime}}) \\ \mathrm{H(u)} \otimes \mathrm {H(u^{\prime})} \Big \downarrow \\ \mathrm {H(C_ {1})} \otimes_ {\mathrm{A}} \mathrm {H(C_ {1} ^ {\prime})} \xrightarrow [ \gamma (\mathrm {C_ {1}} , \mathrm {C_ {1} ^ {\prime}}) ]{} \mathrm {H(C_ {1}} \otimes_ {\mathrm{A}} \mathrm {C_ {1} ^ {\prime}}). \end{array}
\]

设A°为A的反k-代数，并设C°（相应地C'°）为将复形C（相应地C'）视为左（相应地右）A°-模的复形。

\[
\sigma (\mathbf {C}, \mathbf {C} ^ {\prime}): \mathbf {C} \otimes_ {\mathbf {A}} \mathbf {C} ^ {\prime} \rightarrow \mathbf {C} ^ {\prime \circ} \otimes_ {\mathbf {A} ^ {\circ}} \mathbf {C} ^ {\circ}
\]

唯一的次数为0的分次 \(k\) -线性映射，使得对于 \(x \in C_p\) , \(x' \in C_q'\) , \(p, q \in \mathbf{Z}\) ，有

\[
\sigma (\mathbf {C}, \mathbf {C} ^ {\prime}) (x \otimes x ^ {\prime}) = (- 1) ^ {p q} x ^ {\prime} \otimes x.
\]

命题2. — 映射 \(\sigma(\mathbf{C},\mathbf{C}'):\mathbf{C}\otimes_{\mathbf{A}}\mathbf{C}'\to\mathbf{C}^{\prime\circ}\otimes_{\mathbf{A}^{\circ}}\mathbf{C}^{\circ}\) 是 \(k\) -模复形的同构，其逆同构为 \(\sigma(\mathbf{C}^{\prime\circ},\mathbf{C}^{\circ})\) 。

由于映射 \(\sigma(\mathbf{C},\mathbf{C}^{\prime})\) 与 \(\sigma(\mathbf{C}^{\circ},\mathbf{C}^{\circ})\) 互为逆映射，只需证明 \(\sigma(\mathbf{C},\mathbf{C}^{\prime})\) 是复形的态射。现在，对于

\[
x \in \mathrm{C} _ {p} = \mathrm{C} _ {p} ^ {\circ}, \quad x ^ {\prime} \in \mathrm{C} _ {q} ^ {\prime} = \mathrm{C} _ {q} ^ {\prime \circ}, \quad p, q \in \mathbf {Z},
\]

我们有，用 \(D^{c}\) 记 \(C' \otimes_{A^{0}} C\) 的微分，

\[
\begin{array}{l} \sigma (\mathrm{C}, \mathrm{C} ^ {\prime}) \circ \mathrm{D} (x \otimes x ^ {\prime}) = \sigma (\mathrm{C}, \mathrm{C} ^ {\prime}) (d x \otimes x ^ {\prime} + (- 1) ^ {p} x \otimes d ^ {\prime} x ^ {\prime}) = \\ = (- 1) ^ {(p + 1) q} x ^ {\prime} \otimes d x + (- 1) ^ {p + p (q + 1)} d ^ {\prime} x ^ {\prime} \otimes x = (- 1) ^ {p q} d ^ {\prime} x ^ {\prime} \otimes x + (- 1) ^ {p q + q} x ^ {\prime} \otimes d x \\ = (- 1) ^ {p q} \mathrm{D} ^ {\circ} (x ^ {\prime} \otimes x) = \mathrm{D} ^ {\circ} \circ \sigma (\mathrm{C}, \mathrm{C} ^ {\prime}) (x \otimes x ^ {\prime}); \end{array}
\]

这给出 \(\sigma(C, C') \circ D = D^{\circ} \circ \sigma(C, C')\) ，由此得到所期望的断言。

同构 \(\sigma(\mathbf{C},\mathbf{C}^{\prime}):\mathbf{C}\otimes_{\mathbf{A}}\mathbf{C}^{\prime}\to\mathbf{C}^{\circ}\otimes_{\mathbf{A}^{\circ}}\mathbf{C}^{-}\) 称为复形 \(\mathbf{C}\) 与 \(\mathbf{C}^{\prime}\) 的张量积的交换同构。

若 \(u: \mathbf{C} \to \mathbf{C}_{1}\) 与 \(v: \mathbf{C}' \to \mathbf{C}_{1}'\) 是如上所述的两个复形的态射，我们有如下交换图：

\[
\begin{array}{c} \mathrm{C} \otimes_ {\mathrm{A}} \mathrm{C} ^ {\prime} \xrightarrow {\sigma (\mathrm{C} , \mathrm{C} ^ {\prime})} \mathrm{C} ^ {\prime \circ} \otimes_ {\mathrm{A} ^ {\circ}} \mathrm{C} ^ {\circ} \\ u \otimes u ^ {\prime} \Bigg \downarrow \\ \mathrm{C} _ {1} \otimes_ {\mathrm{A}} \mathrm{C} _ {1} ^ {\prime} \xrightarrow [ \sigma (\mathrm{C} _ {1} , \mathrm{C} _ {1} ^ {\prime}) ]{} \mathrm{C} _ {1} ^ {\prime \circ} \otimes_ {\mathrm{A} ^ {\circ}} \mathrm{C} _ {1} ^ {\circ}. \end{array}
\]

假设在本节剩余部分中环 A 是交换的（一般情形参见第 9 节）。

设 C, C′, C′′ 为三个 A-模复形；A-模的典范同态（III，第 64 页）

\[
\varphi : \left(\mathrm{C} \otimes_ {\mathrm{A}} \mathrm{C} ^ {\prime}\right) \otimes_ {\mathrm{A}} \mathrm{C} ^ {\prime \prime} \rightarrow \mathrm{C} \otimes_ {\mathrm{A}} \left(\mathrm{C} ^ {\prime} \otimes_ {\mathrm{A}} \mathrm{C} ^ {\prime \prime}\right)
\]

是复形的同构，这一点可直接由定义验证。

更一般地，设 \((\mathbf{C}^{(i)}, d^{(i)})_{i \in I}\) 是一族A-模复形，其中指标集I是有限且全序的；为记号简便，我们将I等同于N中的区间 \([1, r]\) 。赋予A-模 \(C = \bigotimes_{i=1}^{r} C^{(i)}\) 以如下分次，使得

\[
\mathrm{C} _ {n} = \sum_ {p _ {1} + p _ {2} + \dots + p _ {r} = n} (\mathrm{C} ^ {(1)}) _ {p _ {1}} \otimes (\mathrm{C} ^ {(2)}) _ {p _ {2}} \otimes \dots \otimes (\mathrm{C} ^ {(r)}) _ {p _ {r}},
\]

，并定义C上的一个次数为(-1)的分次A-自同态为

\[
\mathrm{D} \left(x _ {1} \otimes \dots \otimes x _ {r}\right) = \sum_ {j = 1} ^ {r} (- 1) ^ {p _ {1} + \dots + p _ {j - 1}} x _ {1} \otimes \dots \otimes x _ {j - 1} \otimes d _ {j} x _ {j} \otimes x _ {j + 1} \otimes \dots \otimes x _ {r}
\]

其中 \(x_{i}\in (\mathbf{C}^{(i)})_{p_{i}}\)，\(i = 1,\dots,n\)。则(C, D)是一个A-模复形，称为族 \((\mathbf{C}_{i}, d_{i})\) 的张量积复形。对 [0, r] 中的每一个严格递增序列 \(r_{0}, \ldots, r_{k}\)，其中 \(r_{0} = 0\) , \(r_{k} = r\) ，典范结合性同构

\[
\bigotimes_ {j = 0} ^ {k - 1} \left(\bigotimes_ {i = r _ {j} + 1} ^ {r _ {j + 1}} \mathbf {C} ^ {(i)}\right)\rightarrow \bigotimes_ {i = 1} ^ {r} \mathbf {C} ^ {(i)}
\]

是复形的同构。

我们如上定义次数为 0 的分次同态

\[
\gamma ((\mathbf {C} ^ {(i)})): \underset {i \in \mathbf {I}} {\otimes} \mathrm{H} (\mathbf {C} ^ {(i)}) \to \mathrm{H} \left(\underset {i \in \mathbf {I}} {\otimes} \mathbf {C} ^ {(i)}\right) .
\]

注记。--- 3) 可以定义有限族复形的张量积，而无需赋予指标集全序（X，第 185 页，练习 3）。

\begin{enumerate}
\def\labelenumi{\arabic{enumi})}
\setcounter{enumi}{3}
\tightlist
\item
  假设每个 \(\mathbf{C}^{(i)}\) 被赋予一个与其分次结构相容的分次代数结构，且使得 \(d^{(i)}\) 是反导子（III，第117页）。然后让我们赋予 \(\bigotimes_{i\in I}\mathbf{C}^{(i)}\) 以给定结构的左分次张量积代数结构（III，第49页）。则D是一个反导子。事实上，利用张量积的结合性，可以假设I = \{1, 2\}；然后设 \(p_1\) , \(q_1\) , \(p_2\) , \(q_2 \in \mathbf{Z}\) , \(x_1 \in (\mathbf{C}^{(1)})_{p_1}\) , \(y_1 \in (\mathbf{C}^{(1)})_{q_1}\) , \(x_2 \in (\mathbf{C}^{(2)})_{p_2}\) , \(y_2 \in (\mathbf{C}^{(2)})_{q_2}\) ；我们有
\end{enumerate}

\[
\begin{array}{l} \left(\mathrm{D} (x _ {1} \otimes x _ {2})\right) (y _ {1} \otimes y _ {2}) + (- 1) ^ {p _ {1} + p _ {2}} (x _ {1} \otimes x _ {2}) (\mathrm{D} (y _ {1} \otimes y _ {2})) = \\ = (d x _ {1} \otimes x _ {2} + (- 1) ^ {p _ {1}} x _ {1} \otimes d x _ {2}) (y _ {1} \otimes y _ {2}) + \\ \quad + (- 1) ^ {p _ {1} + p _ {2}} (x _ {1} \otimes x _ {2}) (d y _ {1} \otimes y _ {2} + (- 1) ^ {q _ {1}} y _ {1} \otimes d y _ {2}) = \\ = (- 1) ^ {p _ {2} q _ {1}} (d x _ {1}) y _ {1} \otimes x _ {2} y _ {2} + (- 1) ^ {p _ {1} + (p _ {2} - 1) q _ {1}} x _ {1} y _ {1} \otimes (d x _ {2}) y _ {2} + \\ \quad + (- 1) ^ {p _ {1} + p _ {2} + p _ {2} (q _ {1} - 1)} x _ {1} d y _ {1} \otimes x _ {2} y _ {2} + (- 1) ^ {p _ {1} + p _ {2} + q _ {1} + p _ {2} q _ {1}} x _ {1} y _ {1} \otimes x _ {2} d y _ {2} \\ = (- 1) ^ {p _ {2} q _ {1}} [ (d x _ {1}) y _ {1} + (- 1) ^ {p _ {1}} x _ {1} d y _ {1} ] \otimes x _ {2} y _ {2} + \\ \quad + (- 1) ^ {p _ {1} + q _ {1} + p _ {2} q _ {1}} x _ {1} y _ {1} \otimes ((d x _ {2}) y _ {2} + (- 1) ^ {p _ {2}} x _ {2} d y _ {2}) \\ = (- 1) ^ {p _ {2} q _ {1}} [ d (x _ {1} y _ {1}) \otimes x _ {2} y _ {2} + (- 1) ^ {p _ {1} + q _ {1}} x _ {1} y _ {1} \otimes d (x _ {2} y _ {2}) ] \\ = (- 1) ^ {p _ {2} q _ {1}} \mathrm{D} (x _ {1} y _ {1} \otimes x _ {2} y _ {2}) = \mathrm{D} ((x _ {1} \otimes x _ {2}) (y _ {1} \otimes y _ {2}))  . \end{array}
\]

